\documentclass{ltjsbook}
\usepackage{docmute} 
\usepackage{../../myStyle} 

\begin{document}
\chapter{因子分析法}
\label{H03_FA}
\section{因子分析モデル}
\subsection{単因子モデル}
今からお話しするのは，\keyterm{因子分析}{いんしぶんせき}{Factor Analysis}というモデルです。因子分析モデルは古典的テスト理論の拡張であり，もっとも簡単な1因子モデルは次のように表されます。

\[ z_{ij} = a_j f_i + e_{ij} \]

ここで$Z_{ij}$は個人$i$の項目$j$に対する反応を標準得点で表したもの\footnote{\keyterm{標準得点}{ひょうじゅんとくてん}{Standard Score}とは，素点$X_j$を$Z_j = \frac{X_j - \bar{X}}{\sigma_j}$と変換したものであることを思い出してください。標準化されたスコアは平均が0,分散が1になりますので，単位の異なるもの同士であっても標準得点を使うと相互に比較可能になるのでした。}，$a_j$は項目$j$の\keyterm{因子負荷量}{いんしふかりょう}{factor loading}，$f_i$は個人$i$の\keyterm{因子得点}{いんしとくてん}{factor score}，$e_{ij}$は個人$i$と項目$j$の組み合わさった時に生じた誤差と呼ばれます。

\keyterm{因子負荷量}{いんしふかりょう}{factor loading}とは，因子というこのテストで測定したい特性と，項目との関係の強さを表しているものです。
\keyterm{因子得点}{いんしとくてん}{factor score}とは，因子というこのテストで測定したい特性と，個人との関係の強さを表しているもので，その人のスコアだということができます。

記号についている添字に注目してください。添字$i$は個人を，添字$j$は項目を表していますが，因子負荷量は$a_j$と表されています。つまり項目によって変わる変量です。因子得点は$f_i$と表されています。つまり人によって変わる変量です。古典的テスト理論をこの添字を使って表現するならば，$X_{ij} = t_i + e_{ij}$となりますが，これと比べてみると$t_i$が$a_jf_i$に変わったのが因子分析モデルだということになります。$t_i$は個人についての真のスコアなのですが，古典的テスト理論の場合，テストの点数は個人のこの能力だけを反映していると考えられていたことになります。もしテストの問題が難しすぎて，まったく答えることができなければ，その人の能力はゼロということになるわけです。しかし中には悪い問題というのもあって，たとえば小学生に高校で習う知識が必要な問題を解かせるような問題があれば，誰だって解けないかもしれません。解けない問題を出して「学力が低い」と結論づけるのはやや暴力的ですらありますね。このように，古典的テスト理論は項目の良し悪しといったものが評価できないモデルだったのです。

因子分析モデルはこれを改良し，$a_j f_i$としました。すなわち，ある項目に対する反応$z_{ij}$は，その項目が測定したい特徴を十分に反映しているかどうか($a_j$)と，その人がその特徴を有しているかどうか($f_i$)の両方が成立している必要があるわけです。掛け算ですから，一方がゼロであれば結果もゼロになります。すなわち測定したい特徴を反映していない項目($a_j=0$)であれば，どれほどそれについての能力($f_i$)が高くても反応できないのです。たとえば美的センスに非常に秀でた人がいても，数学のテストでその能力を反映させることはできませんよね。これは数学のテストというのが数学力を測定するものであって，美的センスを測定するものではないからです。

因子分析は知能検査や性格検査の文脈から生まれてきたものです。心理学において「知能」とは，何にでも応用可能な一般的な知能というのがあるのか，あるいは語彙力や計算力といった複数の個別の能力があるのか，という議論がありました。知能検査としていろいろなものが考えられますが，それらがきちんと当該能力を測定する検査法だったかどうかもわからないわけで，因子分析モデルはそこを評価できるようにした，ともいえます。性格検査についても同様で，特性論的に考えるならば人間の性質というのは複数のもの，たとえば外向性，神経症傾向，開放性，協調性，勤勉性\footnote{\textcite{Oshio2020}のビッグファイブについての解説に基づいています。}などがあり，ある性格検査の項目は協調性を測定するのには向いているけれども，神経症傾向を測定するには向いていないということがあるわけです。このように，心理学と因子分析モデルは深い関係があります。

\subsection{多因子モデル}
さて先ほどは一因子，あるいは単因子ともいいますが，測定したい特徴が1つだけのモデルでした。学力テストなどは一因子で問題ありません。国語のテストは国語の能力を，数学のテストは数学の能力を測定していれば良いのであって，数学のテストを解くのに美的センス(真美的能力)が必要というのは，むしろ困った状況ですらありますね。しかし，知能検査や性格検査の場合はそうではありません。ある行動傾向，ある形容詞，ある検査がたった1つの能力・性質・心理的要因だけを反映しているとは限りません。たとえば人に優しく振る舞うといっても，その背後に外向性があるのか，あるいはそうすると自分がよく見られるからという利己的な性格があるのか等々が考えられます。1つの項目に複数の要素(因子)が複合的に影響していることを考えるべきです。そこで因子の数が1つではなく，複数ある多因子モデルを考えることにします。多因子モデルは次のように表現されます。

\begin{equation}
	z_{ij} = a_{j1}f_{i1} + a_{j2}f_{i2}+a_{j3}f_{i3} + \cdots + a_{jm}f_{im} + d_j u_{ij}
	\label{factor}
\end{equation}

記号の意味は単因子モデルと基本的には同じです。$z_{ij}$は個人$i$の項目$j$に対する反応を標準得点で表したものであり，$a_{jm}$は第m因子の\keytermL{因子負荷量}{いんしふかりょう}，$f_{im}$は第$m$因子の\keytermL{因子得点}{いんしとくてん}を表しています。因子の数が複数あるモデルですから，$a_{j1},a_{j2},\cdots, a_{jm}$とか$f_{i1},f_{i2},\cdots, f_{im}$のように因子の番号と項目・個人の添字の組み合わせになっていることを確認してください。最後の$e_{ij}$が$d_j u_{ij}$となっていますが，これは誤差についても他の因子と形を同じくし，項目に依存するものとそれ以外に区別しているだけです。

さて，ここでは因子を$m$個あるとしています。この因子はどの項目にも共通して働くので，\keyterm{共通因子}{きょうつういんし}{common factor}と呼ばれます。共通因子がいくつあるかは事前にわかりませんが，一般的にその数は数個〜十数個になります。性格心理学は長い研究の中で，性格を表す言葉に共通する因子はおおよそ5つぐらいであろう，という答えを得るに至りましたが，それ以外の領域では領域ごとの見解があるでしょう。知能が何種類の因子に分かれるのか，あるいはとある心の状態がどういう構造をしているのか，というのは心理学的にみても十分興味のある考え方です。もちろん因子分析によって得られる因子が，人間の潜在的な知能や概念に直接対応しているとはいえないのですが\footnote{むしろテスト項目や調査票などに対する反応パターンが因子として出てくるだけで，性格や知能が数次元あるというより，我々は性格や知能を数次元で捉えることしかできない，という方が正しいでしょう。詳しくは第\ref{H06_Attitude}章で議論します。}，それでも因子がどのような構造(しくみ)をしているのかについての一定の情報を与えてくれます。多因子モデルが心理学一般で広まったのは，こうした心の「構造」に注目する学問との相性が良かったということでしょう。

\section{因子分析の定理}
\subsection{因子分析モデルの展開}
因子分析モデルも古典的テスト理論のように，式の展開から何が見えてくるか考えてみましょう\footnote{以下このセクションは\textcite{kosugi2018}の原稿を再構成したものです。}。

左辺の$z_{ij}$は観測されたデータから算出できるものですが，右辺の因子負荷量，因子得点はいずれも未知数です。データに対して未知数が多すぎるようで，これではどのようにして答えを見つけ出せば良いのかわからないかもしれません。たとえばある人のある項目に対する標準得点が0.12であるとして，それが$0.4 \times 0.3$で得られるのか，$0.2 \times 0.6$で得られるのか，はたまた他の数値の組み合わせで得られるのか，を解く数学的技術は存在しません。この方程式はこのままでは解けないのです。

そこで，この未知数だらけの方程式を解くために，因子について以下のような条件を置きます。
\begin{itemize}
	\item 共通因子の因子得点，独自因子の因子得点は，標準化されている。すなわち，いずれの因子得点も平均点は0であり，分散は1である。
	\item 独自因子は共通因子，他の独自因子と相関しない。
\end{itemize}
この他に，状況に応じて因子同士の間に相関を仮定します。
\begin{itemize}
	\item 共通因子同士の相関を認めないのを「直交因子モデル」，認めるのを「斜交因子モデル」と呼ぶ。
\end{itemize}

このような仮定を置いたら問題が解けるようになるのでしょうか？ 実はこの問題を解く鍵は，多変量データであればなんとかなるのです！

2つの変数，$j$と$k$の標準得点から，
\begin{equation}
	r_{jk}=\frac{1}{N} \sum_{i=1}^{N} z_{ij}z_{ik}
\end{equation}

のように，相関係数が算出されることを思い出してください。先ほどの因子分析の基本代数式(式\ref{factor})をこの式に代入してみましょう。
\begin{equation}
	\begin{array}{ccl}
		r_{jk} & = & \displaystyle{\frac{1}{N} \sum_{i=1}^{N} z_{ij}z_{ik}}                                                                                                                  \\
		       & = & \displaystyle{\frac{1}{N} \sum_{i=1}^{N}}(a_{j1}f_{i1}+a_{j2}f_{i2}+ \ldots \ +a_{jm}f_{im}+d_{j}u_{ij})(a_{k1}f_{i1}+a_{k2}f_{i2}+ \ldots \ +a_{km}f_{im}+d_{k}u_{ik})
	\end{array}
\end{equation}

これは代数の計算としてやっていくと，非常に煩雑で間違いが起きやすそうです。そこで，少し視覚化してわかりやすくしてみましょう。
多項式の掛け算は，各項目の総当たり戦ですので，列方向に$z_{ij}$，行方向に$z_{ik}$の各要素を置いて，要素同士の組み合わせ表を作ります(図\ref{x03_multiply})。

\begin{figure}[ht]
	\[
		\begin{array}{rl|cccccc|}
			        &                                                        & \multicolumn{6}{c}{z_{ij}=}                                                                                                                                                                                                                                        \\
			        &                                                        & a_{j1}f_{i1} +                                          & a_{j2}f_{i2}+                                            & a_{j3}f_{i3} +                                         & \cdots\cdots +                                          & a_{jm}f_{im} + & d_jU_{ij} \\\hline
			        & \begin{array}{c}a_{k1}f_{i1}\\+\end{array}             & \multicolumn{1}{>{\columncolor{red!9}}c}{}              & \multicolumn{4}{>{\columncolor{cyan!9}}c}{}              & \multicolumn{1}{>{\columncolor{green!9}}c}{}                                                                                                  \\
			z_{ik}= & \begin{array}{c}a_{k2}f_{i2}\\+\end{array}             & \multicolumn{1}{>{\columncolor{cyan!9}}c}{}             & \multicolumn{1}{>{\columncolor{red!9}}c}{\circled{1}}    & \multicolumn{3}{>{\columncolor{cyan!9}}c}{\circled{2}} & \multicolumn{1}{>{\columncolor{green!9}}c}{}                                         \\
			        & \begin{array}{c}a_{k3}f_{i3}\\+\end{array}             & \multicolumn{2}{>{\columncolor{cyan!9}}c}{}             & \multicolumn{1}{>{\columncolor{red!9}}c}{}               & \multicolumn{2}{>{\columncolor{cyan!9}}c}{}            & \multicolumn{1}{>{\columncolor{green!9}}c}{\circled{3}}                              \\
			        & \hspace{10pt}\begin{array}{c}\vdots\\+\end{array}      & \multicolumn{3}{>{\columncolor{cyan!9}}c}{\circled{2}}  & \multicolumn{1}{>{\columncolor{red!9}}c}{}               & \multicolumn{1}{>{\columncolor{cyan!9}}c}{}            & \multicolumn{1}{>{\columncolor{green!9}}c}{}                                         \\
			        & \begin{array}{c}a_{km}f_{im}\\+\end{array}             & \multicolumn{4}{>{\columncolor{cyan!9}}c}{}             & \multicolumn{1}{>{\columncolor{red!9}}c}{}               & \multicolumn{1}{>{\columncolor{green!9}}c}{}                                                                                                  \\
			        & \hspace{7pt}\begin{array}{c}\\ d_kU_{ik}\\ \end{array} & \multicolumn{5}{>{\columncolor{green!9}}c}{\circled{3}} & \multicolumn{1}{>{\columncolor{yellow!9}}c}{\circled{4}}                                                                                                                                                 \\\hline
		\end{array}
	\]
	\caption{項目同士の総当たりを考える}
	\label{x03_multiply}
\end{figure}

図\ref{x03_multiply}に示されたのは個人$i$についてのものであり，これが人数分ある，すなわち$\sum_{i=1}^N$をつけないといけないことに注意してください。
さて図を軽く色分けしてあるのですが，ここにあるように計算すべき領域を4つに分けて考えていきましょう。
\begin{description}
	\item[\circled{1}の領域] 因子$p$と$p$の積和部分(同じ因子同士の掛け合わせ)
	\item[\circled{2}の領域] 因子$p$と$q$の積和部分(異なる因子同士の掛け合わせ)
	\item[\circled{3}の領域] 因子$p$($q$)と独自因子の積和部分
	\item[\circled{4}の領域] 独自因子同士の積和部分
\end{description}

この各パートを順に計算していきましょう。まず\circled{1}ですが，たとえば第一因子については
\begin{equation}
	\frac{1}{N}\sum_{i=1}^{N}a_{j1}a_{k1}F_{i1}^2
\end{equation}
となることがわかります。ここで，$a_{j1}$と$a_{k1}$は$N$には関係がない($\sum$は$i$が$1$から$N$まで変化することを意味しているが，係数$i$はどちらにも入っていない)ので，総和して割る意味がないことに気づきます。そうなると，必然的にこの式は，
\begin{equation}
	\frac{1}{N}\sum_{i=1}^{N}a_{j1}a_{k1}F_{i1}^2=a_{j1}a_{k1}\frac{1}{N}\sum F_{i1}^2
\end{equation}
となります。この$F_{i1}$は因子得点であり，上の仮定より標準化されたものだということになります。さて，標準得点と標準得点の積和平均は相関係数になることをもう一度思い出してください！そうすると，これは自分自身との相関係数を表していることになりますから，当然$F_{i1}^2=1.0$であることがわかります。

結局，\circled{1}のエリアは
\begin{equation}
	\frac{1}{N}\sum_{i=1}^{N}a_{j1}a_{k1}F_{i1}^2=a_{j1}a_{k1}\frac{1}{N}\sum F_{i1}^2=a_{j1}a_{k1}
\end{equation}
と，とてもあっさり書き下すことができるのです。

つづいて\circled{2}を見てみましょう。ここは異なる因子がかけ合わさった部分ですね。落ち着いて，第一因子と第二因子を例にして考えてみましょう。この箇所で得られるのは，
\begin{equation}
	\frac{1}{N}\sum a_{j1}F_{i1} a_{k2}F_{i2} = a_{j1}a_{k2}\frac{1}{N}\sum F_{i1}F_{i2}
\end{equation}
ということになります。ここで，$F_{i1}$および$F_{i2}$はそれぞれ第一，第二因子における個人$i$の因子得点を意味しています。因子得点は標準化されていることをもう一度思い出すと，これは第一因子と第二因子の相関係数になります。ここで，この因子分析が直交因子モデルだと考えますと，因子同士に相関がないわけですから，数字としては$0.0$で消えてしまいます。するとこの部分は，

\begin{equation}
	\frac{1}{N}\sum a_{j1}F_{i1} a_{k2}F_{i2} = a_{j1}a_{k2}\frac{1}{N}\sum F_{i1}F_{i2}=0
\end{equation}
となることがわかりました。つまり，この領域\circled{2}は，すべて$0$になってしまうのです。


続いて\circled{3}の部分について考えてみましょう。これはある共通因子と独自因子の積和部分です。例によって標準得点同士の関係から，相関係数を算出することになりますが，独自因子は共通因子と無相関であることを考えると，
\begin{equation}
	\frac{1}{N}\sum a_{i1}F_{i1}d_{j}U_{ij}=a_{ik}d_{j} \frac{1}{N}\sum U_{ij}F_{i1} =0
\end{equation}
とこのように，この箇所もすべて$0$になってしまいます。

最後の\circled{4}に至っては，独自因子と独自因子の積和ですから，これも
\begin{equation}
	\frac{1}{N}\sum d_{j}d_{k}U_{ij}U_{ik}=d_{j}d_{k}\frac{1}{N}\sum U_{ij}U_{ik}=0
\end{equation}
のように$0$になります。結局，消えて無くなるのがほとんどで，残るのは\textcircled{1}の部分だけであり，$r_{jk}$を考えるときはそこだけ考慮すればよいことになります。

整理すると，
\begin{equation}
	r_{jk}=a_{j1}a_{k1}+a_{j2}a_{k2}+ \cdots +a_{jm}a_{km}
\end{equation}
ということがわかります。つまり，項目$j$と項目$k$の相関係数は，項目$j$の因子負荷量と項目$k$の因子負荷量を，すべての因子について総和したものであるということです。因子分析の基本モデルから導出されるこの定理を，とくに\textbf{因子分析の第二定理}と呼びます。

ここで同じ項目同士の相関を考えてみましょう。項目$j$と項目$j$の相関係数は，もちろん$1.0$になりますね。これを因子分析の基本式で表すと，次のように表現できます。
\begin{equation}
	\label{facteo1}
	r_{jj}=a_{j1}^2+a_{j2}^2+ \cdots + a_{jm}^2 + d_j^2 = 1.0
\end{equation}

さて，この式が意味するのはなんでしょうか。意味を考えてみると，ある項目それ自身との相関係数は，因子負荷の二乗和からなっている，ということがわかります。これこそ\textbf{因子分析の第一定理}と呼ばれるものであり，解けるはずのなかった方程式を解くための鍵となる式なのです。

\section{因子分析の定理}
\label{FactorAnalysisTheorems}
数式の展開はいったんここまでにして，
第一定理は次のような形をしているのでした。
\[a_{j1}^2+a_{j2}^2+ \cdots + a_{jm}^2 + d_j^2 = 1.0 \]
ここで共通因子部分を，$a_{j1}^2+a_{j2}^2+ \cdots + a_{jm}^2=h_j^2$のようにすると，この式は単に$h_j^2 + d_j^2 = 1.0$となります。この$h_j^2$のことをとくに\keyterm{共通性}{きょうつうせい}{communality}といいますが，この式は共通性と独自因子の二乗和が$1.0$になることを意味しています。言い換えると，全体を100\%とした比率で共通性と誤差を比較できるということです。
共通性は因子負荷量の二乗和で，共通因子はそのテストの背後にある共通の要因，すなわちテストで測定したいものだったわけです。古典的テスト理論では，モデル式の展開から信頼性を全分散中に示る真のスコアの割合と定義しましたが，因子分析モデルはこのように1つの項目$j$における共通因子の割合を算出し，項目の信頼性を考えることができます。因子分析モデルにおける信頼性は，1項目の中に含まれる共通因子の大きさだともいえるわけです。逆に$d_j^2=1-h_j^2$は\keyterm{独自性}{どくじせい}{uniqueness}は，当該項目がそのテストで測っていないものの大きさを表しており，この割合があまりにも大きいと「この項目は全然関係ないものを測っちゃってるんじゃないか？ 」と疑われることになります。多因子モデルにおいては，多角的に対象を切り分けるために多くの質問を投げかけるわけですが，独自性の高い項目は回答者に負担をかけるだけの邪魔なものですから，実践上はこうした項目を除外することが少なくありません。因子分析には\keyterm{単純構造の原則}{たんじゅんこうぞうのげんそく}{principle of simple structure}と呼ばれるものがあり，項目は該当する因子を適切に反映し，かつ，他の因子と関係ないことが美しいとされます。尺度構成段階では，共通性(独自性)をみて項目の良し悪しが判断されるのです。

次に第二定理を見てみましょう。第二定理は次のような形をしているのでした。
\[  r_{jk}=a_{j1}a_{k1}+a_{j2}a_{k2}+ \cdots +a_{jm}a_{km} \]
2つの項目$j$と$k$の相関係数は，それぞれの因子負荷量の積和の形で表される，というものです。ここに誤差の話は入ってこず，共通因子だけで話ができています。

左辺の相関係数は，2つの項目がどれほど同じものを測定しているかの指標です。相関係数(の絶対値)が高ければ高いほど，2つの項目は同じものを指し示しているわけです。逆に相関係数が低いということは，2つの項目に関係がないことを表します。ここで右辺に目をやりますと，右辺の各項目は因子負荷量の積の形になっています。左辺の値が小さくなる1つの理由は，ある共通因子$m$が項目$j,k$に対して，異なる方向で寄与しているからだと考えられるでしょう。そしてそのパターンが一貫していないという状況です。そもそも相関係数が小さいところからは因子を見つけ出すのは難しいのですが，そうした状況があるのはある項目ペアについて因子同士の向きがバラバラに影響しているからだといえます。そのような状況は，測定がきちんとできているかどうか怪しいですね。\textbf{測定の一義性}ともいわれますが
% \footnote{公理論的測定論における測定論の基礎問題には，表現の問題，一意性の問題，有意味性の問題，尺度構成の問題があり\parencite{Coombs1970}，一意性の問題とは「測定手順があるとして，対象に数値を割り当てる時にどの程度まで自由に行えるか」あるいは「与えた数値の相互変換はどのようなものが許容されるか」といったことを指します。}
，そのような尺度は妥当性が低いといえるでしょう。

このように，因子分析モデルは第一定理で信頼性を，第二定理で妥当性をあらわすものになっているのです。


\section{因子分析の歴史と展開}

ところで，因子分析モデルもテスト理論も，潜在変数モデルとしては同じでいずれも古典的テスト理論からの発展系なのでした。
因子分析モデルは多段階反応が一般的で，多因子モデルで「潜在的な(心理学的な)構造はどうか」ということを問題にします。ここでの目的は全体に共通する要素やその構造であり，何種類に別れて要素間の関係はどうなっているのか，というところが中心的関心事になります。
一方，項目反応理論はバイナリ反応が一般的で，因子数はひとつです。学力テストはその学力が測定できていることが重要で，因子の構造よりも因子得点をより精緻に推定できることの方が重要だからです。因子分析の言葉で表すなら，因子得点をより精緻に表現しようとしているわけです。

さて，GRMは，項目反応理論の多段階モデルでした。実はGRMは因子分析モデルの発展系でもあるのです。
因子分析は相関係数のモデルであったことを思い出してください。因子分析モデル自体は$z_{ij}$から始まっていましたが，変数同士の共分散$r_{ij}$を考えるといくつかの仮定から\keytermL{因子負荷量}{いんしふかりょう}だけのモデルに簡略化され，推定できるようになるのでした\footnote{具体的にどうやって因子負荷量を算出するかは次回以降のお楽しみです。}。この標準化された共分散，すなわち相関係数はピアソンの積率相関係数ともいわれ，間隔尺度水準以上の数値を使って計算されます。多段階の反応は順序尺度水準ですから，相関係数を計算するのは不適切で，そのまま因子分析することはできません\footnote{できないのですが，尺度値に変換することもなく機械的に分析してしまう悪い習慣が蔓延っているのは何度も指摘している通りです。くどいと思われるかもしれませんが，私は憤っているのです。}。では\keytermL{順序尺度水準}{じゅんじょしゃくどすいじゅん}の相関係数は，というとスピアマンの順序相関などを思い出す人がいるかもしれませんが，より統計的なモデルで表現しやすいものがあります。

順序尺度水準の変数$\times$順序尺度水準の変数の相関は\keyterm{ポリコリック相関係数}{ぽりこりっくそうかんけいすう}{polychoric correlation}といいます。
順序尺度水準の変数$\times$間隔尺度水準の変数の相関は\keyterm{ポリシリアル相関係数}{ぽりしりあるそうかんけいすう}{polyserial correlation}といいます。
ついでにバイナリ変数$\times$バイナリ変数の相関係数は\keyterm{テトラコリック相関係数}{てとらこりっくそうかんけいすう}{tetracholic correlation}といいます。

これらの相関係数はいずれも，順序(あるいはバイナリ)変数の背後に連続体があると考え，潜在的な連続体$\times$潜在的な連続体の相関係数を連続体のカテゴリが変わる閾値とともに推定するのです(図\ref{fig::05_05})。
\begin{figure}[ht]
	\begin{center}
		\includegraphics[width=0.8\linewidth,keepaspectratio,clip]{../images/Slide05_01.png}
		\caption{ポリコリック相関係数のイメージ}
		\label{fig::05_05}
	\end{center}
\end{figure}

こうして推定された相関係数を使って因子分析を行うと，順序尺度水準に適した因子分析を行うことができます。この方法をとくに\keyterm{カテゴリカル因子分析}{かてごりかるいんしぶんせき}{categorical factor analysis}といいます。もっともこの名前を覚えておく必要はありません。\textbf{カテゴリカル因子分析は段階反応モデルと数学的に等価である}ことがわかっています。GRMをすることはカテゴリカル因子分析をしていることと同じ，なのです。

もっともGRMは項目反応理論の系列ですから，単因子構造を仮定しています。複数の因子を想定するカテゴリカル因子分析に対応するのは，正確には\keyterm{多次元項目反応理論}{たじげんこうもくはんのうりろん}{Multidimensional Item Response Theory}といいます。しかし数学的・技術的には同じであり，カテゴリカル因子分析をした結果からIRCCCを描くこともできますし，実際に分析するソフト上では使用する変数がカテゴリカル(順序尺度水準)であることを指定するだけです。われわれユーザはもはや悩む必要はなく，ただただデータに適した分析をするだけで良いのです。

\subsection{系譜の違いはどこに関係するか}

因子分析モデルと項目反応理論が，カテゴリカル因子分析として概念的に統合されることを話してきました。本質的にはこのように違いがないのですが，系譜の違い，出自の違いはそれぞれの利用される文脈で，何を強調するかに影響してくることがあります。

たとえば因子分析の文脈では，共通性が低い項目は削除し，綺麗な因子構造を目指そうという考え方があります。尺度作成の中で1つの項目は1つの因子に対応しているべきであるという考え方があり(\keyterm{単純構造の原理}{たんじゅんこうぞうのげんり}{Principle of Simple Structure}といいます)，もし1つの項目が複数の因子の影響を受けているようであれば，「美しくないので」削除されることがあります。因子分析は知能，性格，態度の研究で展開されてきたため，美しい「構造」を見つけ出すことに狙いがありますから，この美しさを損なうもの(項目)は取り除く，という方向に行きがちです。

一方で，項目反応理論はテスト理論の生まれです。もちろん回答者の能力や特性を測定するのに優れた項目とそうでない項目，という峻別はしますが，中でも「この項目は測定能力の偏差値30程度の回答者を測定するのに適している」とか，「偏差値70程度の回答者を測定するのに適している」と判断できます。偏差値30や70を測定するのに適した項目とは，非常に簡単(ほとんどの人が正答する)だったり，非常に難易度が高い(ほとんどの人が誤答する)項目です。心理尺度でいうと床効果，天井効果がみられる項目とされる，どちらかに偏った分布をもつ項目です。しかし，それを捨てるということにはならず，どのような項目でも回答者の能力を推定するための情報量がゼロではない，という考えから，さまざまな項目をどんどんためていく方向にいきがちです。項目反応理論の文脈では，あらゆる人に対する測定を準備しておく必要があり，むしろ回答者の特性にあわせて設問の方を選んだり，事前の項目特性から，前もってテストの項目構成をデザインする，ということを目的とするのです。

このように使われるシーンによって，「構造」か「機能(=得点)」のどちらに注目するかが変わり，結果的に実践の方針がちがってくることもあるのです。図\ref{fig::05_06}にこの心理学的系譜(左ルート)とテスト理論的系譜(右ルート)の流れを描いてみました。

ポイントは「最終的には同じところに辿り着く」という点ですから，歴史的流れや個々のモデルの細かい数式を完全に理解していなくてもいいかもしれません。それよりは，心理学者として，あるいはテストをする側として，回答者に無理のない，それでいて必要な程度に精緻な情報が得られる適切な分析方法を選択できるようになることが重要です。
\begin{figure}[ht]
	\begin{tikzpicture}
		\node[draw,align=center] (a) at (5,0)
		{古典的テスト理論\\$X=t+e$};
		\node(P0) at (5,0.9)[above] {見えないものを測定・数値化したい};
		\node[draw,align=center] (b) at (0,-2.5) {探索的因子分析\\$z_{ij}=\sum a_{jm}f_{im} + d_jU_{ij}$};
		\node[draw,align=center] (c) at (10,-2.5) {項目反応理論\\ $p_j(\theta) = \frac{1}{1+exp(-1.7a_j(\theta_i -b_j))}$};
		\node[draw] (d) at (5,-5) {カテゴリカル因子分析};
		\draw[->] (a) --(b);
		\draw[->] (a) --(c);
		\draw[->] (b) --(d);
		\draw[->] (c) --(d);
		\node (P1) at (1.3,-1.25)[above] {構造に興味がある};
		\node (P2) at (8.7,-1.25)[above] {機能に興味がある};
		\node (P3) at (2.5,-4.2)[left,align=center] {順序尺度水準へ\\得点の精緻化};
		\node (P4) at (7.5,-4.2)[right,align=center] {多段階反応へ\\多次元に拡張};
	\end{tikzpicture}
	\caption{理論・モデルの流れ。左側のルートが心理学的系譜，右側のルートがテスト理論的系譜}
	\label{fig::05_06}
\end{figure}

\section{調査研究の手順}
最後に，一般的におこなわれている心理尺度作成の手順を大まかに理解しておきましょう。

\paragraph{構成概念の設定}心理尺度を作り始めるにあたって，まずは何を測りたいのかを明確にする必要があります。\keyterm{構成概念妥当性}{こうせいがいねんだとうせい}{Construct Varidity}についてのそもそも論ですね。ある概念を測定したいとして，そういう概念が本当に存在するのか，どこの誰がどのように持つ概念なのかを明確にする必要があります。たとえば文部科学省がスローガンのように掲げる\textbf{生きる力}というのを測ってみたいと思っても，それは何なのかわかりません。死んでなければ生きているのですから，生きる力はありそうに思えます。でもそういうことじゃないらしい。じゃあどういうことなのか，ということを考えて行かなければなりません。あるいは，みなさんは\textbf{いちびり}という関西弁をご存知ですか。いたずらっ子，わざと変なことをして注目を引きたい子，のような意味なのですが，このいちびり感は関西の人間にしかわからない感覚かもしれません。であれば調査対象者が限定的な，一般的ではない概念ですから，因子分析の行う多くの人の反応から共通成分を抜き出すという考え方には適さない概念です。より一般的なものに考え直す必要があるでしょう。

ほかにも，心理学的な態度なのか，パーソナリティのような行動の傾向なのか，抑うつ傾向のように心身両方の問題なのか，そしてそれらが紙とペンで測定可能なものなのかどうか，改めて考えましょう。何を当たり前のことを，と思うかもしれませんが，意外とおかしな問題設定に入り込んでないとも限らないのです。また，心理尺度や心理学的な概念はすでにさまざまなものが存在します。これらを見比べて，自分の測定したい概念が他の概念とどのように同じで，どのように違うのかを論理的に考えられなければなりません。どこかにある尺度とほとんど同じであればオリジナリティが存在しないだけでなく，車輪の再発明というムダな努力をすることになります。完全にオリジナルなものを考えるのは難しいかもしれませんが，新しい概念を使うことによってこれまでの概念で説明できたことを含み，さらにこれまでの概念で説明できなかったことも説明できるようになる，という利点が必要です。これらは\keyterm{弁別的妥当性}{べんべつてきだとうせい}{distincive validity}にも関わる問題です。

\paragraph{項目の選出}測定したい概念が明確になれば，これに関わる項目を作り出す必要があります。ある態度を強く持っている人はどのような振る舞いをし，どのような振る舞いをしないのか。どのような意見に賛成し，どのような意見には反対するのか，といったことをさまざまな角度から検証します。「目に見えないものを測定する」のが心理測定であり，質問紙調査というのはこの見えない的に向かって項目という矢を射掛けるようなものです。矢の数はなるべく多く，また人は嘘をついたり間違えたりする生き物ですから多角的に聞くことで，捉えるべき本質に迫っていく必要があります。言葉を使って聞くわけですから，古すぎる・新しすぎる表現は避け，専門用語は使わずなるべく平易な言葉で聞く必要があります。ワーディング(言い回し)についても細部まで注意しましょう\footnote{たとえば「テレビやラジををよく見聞きしますか」という質問は悪い質問です。テレビしか見ない人，ラジオしか聞かない人がどう答えていいか悩むからです。これはダブルバーレルと呼ばれる悪手ですが，ほかにも色々ありますので調査法の専門書を参考にしてみてください。}。

\paragraph{予備調査の実施}ある程度項目が集まれば，予備調査を行います。用意した項目の5倍から10倍の人を集めるのが良い，と一般的にいわれています\parencite{kosugi2016a}。ここでの予備調査項目は，十分考えられたものではあると思いますが，分析してみると意外と不適切な項目というのが出てくるものですので，多めに用意されていることでしょう。必然的に，予備調査の対象サイズも百人以上の大きなものになるのが一般的です。

\paragraph{探索的因子分析}ここが尺度作成に関わる分析のメインです。ここで因子数を決める(あるいは確認する)ことになり，因子的な妥当性みたり，不適切な項目は除外したりします。別の聞き方の方が良いような項目があれば，項目の入れ替えを行なって再調査することもあります。この探索的な因子分析と予備調査を繰り返して，何度行っても同じ因子構造になり，同じ項目が同じ因子に含まれるというのが確認できれば，項目セットは完成するといえるでしょう。

\paragraph{項目反応理論}ある因子に関連する項目とそのデータセットを抜き出して，項目反応理論(段階反応モデル)を実行しましょう。なぜ一部を抜き出すかというと，項目反応理論モデルは1次元性を仮定したモデルであることが基本だからです。これを多次元に展開した多次元段階反応モデルもありますので，直接そちらを使っても構いません。ともかくこれを行うことで，反応段階数を確認でき，また項目情報曲線やテスト上表曲線を書くことで，この尺度がどのような領域を得意とする項目群からなるのかを記述できることになります。

\paragraph{本調査へ}最終的に項目群が確定したら，大規模な調査を行ってテストの\keytermL{標準化}{ひょうじゅんか}を目指します。ここでいう標準化とは，標準得点にする数値的演算ではなく，大規模調査に基づいてこの尺度がどのような得点分布をするかに基づき，ロウデータを採点結果としての点数に換算する手続きを明確化する，という意味です。因子構造などはほぼ確定しているはずですから，このテストを使うとどのようなデータの散らばりが得られるのかを確定することができるはずです。この最後のステップの目的はこのテストで測定される人の平均や散らばり，分布の形状を確認し，確定することにあります。使うときに逐一分析しなくてもいいように，項目回答パターンがどれぐらいであれば，全体のどのあたりに位置するかといったプロフィールをつくるのです。膨大なデータに裏付けられた採点なので，尺度の値が何らかの目安になり得るわけです。

以上が大体の流れになります。ここにあるように，予備調査を何度も繰り返して因子構造を確定し，そこから本調査を行うことで，ほぼこの尺度はどれぐらいの点数で分布するか，といったことがわかるようになります。
かつては因子分析を1回実行するだけでも多大な時間がかかりましたから，心理尺度を作るというのはそれだけでライフワークになるような作業でした。幸い最近は分析のスピードが飛躍的に向上しましたので，簡単に分析してやり直すということができます。しかしだからといって，構成概念妥当性を疎かにしてはいけませんし，「やったらこうなった」というようなやり逃げ研究になるのではなく，使うための尺度を作るのだということは忘れてはいけないでしょう。

\end{document}